\section{Approach}\label{sec:approach}
\subsection{Overview}\label{sec:approach-overview}
Figure~\ref{fig:approach-overview} organizes \tool{} into three tiers. An access tier separates an AI agent, which reaches the system only through restricted MCP tools, from the operator, who executes probes, reviews proposals and requests dispatch. The middle tier is the \tool{} readiness layer, whose six numbered components connect an experiment plan to an admission decision. The bottom tier is the inherited execution substrate: the pinned \upstream{} templates and the Slurm scheduler of the shared cluster. \tool{} adds the readiness layer and reuses the substrate unchanged.

\begin{figure}[t]
\centering
% RES-AI architecture overview (Figure fig:approach-overview); requires tikz.
\begin{tikzpicture}[
  font=\tiny, >={Stealth[length=3.5pt,width=2.6pt]}, line width=0.45pt,
  comp/.style={draw, rounded corners=1pt, fill=white, text width=2.05cm,
    minimum height=1.5cm, align=center, inner sep=2pt},
  num/.style={circle, fill=black, text=white, inner sep=0pt, minimum size=7.5pt,
    font=\tiny\bfseries},
  term/.style={draw, rounded corners=4pt, fill=gray!8, align=center, inner sep=2.5pt},
  inherited/.style={draw, fill=gray!15, align=center, inner sep=2.5pt},
  actor/.style={draw, fill=white, align=left, inner sep=2.5pt},
  flow/.style={->},
  rev/.style={->, dashed},
  lbl/.style={fill=white, inner sep=0.8pt},
]
% Readiness layer: (1) plan -> (3) probes -> (4) ledger -> (5) assessment -> (6) dispatch.
\node[comp] (plan) at (1.05,0) {\textbf{Experiment plan}\\[2pt]
  $P=(S,X,O,B,I)$\\ science $S$ apart\\ from execution $X$\\ identity $h=H(P)$};
\node[comp] (smoke) at (3.65,0) {\textbf{Bounded smoke}\\ \textbf{probes}\\[2pt]
  run obligations $O$\\ in time-limited\\ allocations};
\node[comp] (ledger) at (6.25,0) {\textbf{Evidence ledger}\\[2pt]
  records $e_i$ with\\ $h_i,t_i,C_i,m_i,c_i$;\\ failed and obsolete\\ attempts kept};
\node[comp] (assess) at (8.85,0) {\textbf{Readiness}\\ \textbf{assessment}\\[2pt]
  Eq.~(\ref{eq:ready}): identity,\\ coverage, outputs,\\ freshness, memory,\\ stability, budgets};
\node[comp] (dispatch) at (11.45,0) {\textbf{Guarded dispatch}\\[2pt]
  reassess $(P,E)$;\\ verify inputs,\\ graph, resources;\\ dry run by default};
\node[comp, minimum height=1.0cm] (discover) at (1.05,-2.05) {\textbf{Resource}\\ \textbf{discovery}\\[2pt]
  node, partition and\\ topology reads;\\ snapshot $G_t$};
\node[comp, dashed, minimum height=0.8cm, text width=2.6cm] (refine) at (6.25,1.8)
  {\textbf{Refinement proposal}\\[1pt] OOM and declared\\ batch invariance\\ $\Rightarrow$ halve batch in $X$;\\ $S$ unchanged};
\node[term] (blocked) at (9.9,-2.05) {\textbf{BLOCKED}\\ with reasons};

\foreach \n/\k in {plan/1, discover/2, smoke/3, ledger/4, assess/5, dispatch/6}
  \node[num] at (\n.north west) {\k};

\draw[flow] (plan) -- node[lbl, above=1pt] {$P,h$} (smoke);
\draw[flow] (smoke) -- node[lbl, above=1pt] {$e_i$} (ledger);
\draw[flow] (ledger) -- node[lbl, above=1pt] {$E$} (assess);
\draw[flow] (assess) -- node[lbl, above=1pt] {\textsc{ready}} (dispatch);
\draw[flow] (discover) -- node[lbl, right=1pt] {$G_t$ informs reviewed $X$} (plan);
\draw[flow] (assess.south) |- (blocked.west);
\draw[flow] ([xshift=-0.5cm]dispatch.south) |- node[lbl, pos=0.3, left=1pt] {check fails} (blocked.east);
\draw[rev] (assess.north) |- node[lbl, pos=0.75, above=1pt] {\textsc{revalidate}} (refine.east);
\draw[rev] (refine.west) -| node[lbl, pos=0.25, above=1pt, align=center]
  {$P'$ after operator review;\\ new $h$, fresh evidence} (plan.north);

\node[fit=(plan)(dispatch)(discover)(blocked)(refine), draw, rounded corners=3pt,
  inner xsep=6pt, inner ysep=7pt, line width=0.7pt] (layer) {};
\node[anchor=north east, inner sep=3.5pt] at (layer.north east) {\tool{} readiness layer};

% Access tier: agents observe and propose; the operator executes and submits.
\node[actor, anchor=south west] (agent) at ([yshift=0.6cm]layer.north west)
  {\textbf{AI agent} (MCP client)\\ tools \code{snapshot}, \code{assess}, \code{propose}; no execution};
\node[actor, anchor=south east] (operator) at ([yshift=0.6cm]layer.north east)
  {\textbf{Operator}\\ runs probes, reviews proposals, requests dispatch};
\draw[->, densely dotted] (agent.south) -- node[lbl, right=1pt] {observe / propose}
  (agent.south |- layer.north);
\draw[->] (operator.south) -- node[lbl, left=1pt] {execute / approve}
  (operator.south |- layer.north);

% Inherited execution substrate.
\node[inherited, anchor=north east, minimum width=9.8cm] (usehpc)
  at ([yshift=-0.6cm]layer.south east)
  {\upstream{} execution templates (inherited, pinned)\\
   job-graph expansion $\cdot$ Slurm runner $\cdot$ audit records $\cdot$ submitter};
\node[inherited, fill=gray!30, anchor=north, minimum width=12.6cm] (slurm)
  at ([yshift=-0.45cm]layer.south |- usehpc.south)
  {Slurm scheduler on a shared GPU cluster; allocation and admission remain under site policy};
\draw[flow] (smoke.south) -- node[lbl, pos=0.55, right=1pt] {probe jobs} (smoke.south |- usehpc.north);
\draw[flow] ([xshift=0.5cm]dispatch.south) -- node[lbl, pos=0.9, left=1pt] {explicit submit}
  ([xshift=0.5cm]dispatch.south |- usehpc.north);
\draw[<->] (discover.south) -- node[lbl, pos=0.4, right=1pt] {read-only queries} (discover.south |- slurm.north);
\draw[flow] (usehpc.south) -- node[lbl, right=1pt] {\code{sbatch}} (usehpc.south |- slurm.north);
\end{tikzpicture}

\caption{Layered architecture of \tool{}. Numbered components form the readiness layer; shaded tiers are inherited and unchanged. Solid arrows carry data or decisions: the plan $P$ and its identity $h$, attempt records $e_i$, accumulated evidence $E$ and the resource snapshot $G_t$. Dashed elements form the only revision path, which requires operator review, yields a new identity and invalidates earlier evidence for admission. The dotted agent connection permits observation and proposals but neither execution nor submission. OOM denotes out of memory.}\label{fig:approach-overview}
\end{figure}

The experiment plan (1) separates scientific choices $S$ from execution settings $X$ and records the mandatory probes $O$, validation policy $B$ and input identities $I$; its hash $h$ identifies the exact plan. Resource discovery (2) issues read-only scheduler queries and archives a snapshot $G_t$ that informs the reviewed choice of $X$. The scheduler retains control of allocation and admission.

Bounded smoke probes (3) exercise the declared obligations through the \upstream{} runner under explicit, time-limited allocations started by the operator. Each attempt yields a record $e_i$ in the evidence ledger (4), carrying its plan identity, timestamp, passed obligations, output validity, peak memory fraction and cost. Failed and obsolete attempts remain in the ledger for validation-cost accounting. Readiness assessment (5) evaluates Equation~\eqref{eq:ready} over the current plan and accumulated evidence $E$: identities must match, the two most recent observations must be fresh, complete, valid and stable, and attempt and compute budgets must hold. Missing or invalid evidence yields BLOCKED with reasons.

A READY decision passes to guarded dispatch (6), which reassesses readiness immediately before dispatch and verifies declared files, the resolved execution graph and the supported resource mapping. Any failed check yields BLOCKED. Otherwise the wrapper produces a dry run by default and invokes the upstream submitter only on explicit operator request.

The dashed path applies only when a permitted refinement exists. After a matching out-of-memory failure, assessment returns REVALIDATE with a proposal that halves the batch size in $X$, provided the scientific contract declares batch invariance; $S$ is never altered. The operator reviews the proposal and adopts the revised plan $P'$, whose new identity prevents reuse of earlier evidence, and collects fresh matching observations. Without an applicable proposal, the operator must diagnose the failure, obtain missing evidence or stop. Agents may call discovery, assessment and proposal tools, whereas probe execution and submission remain operator actions. Algorithms~\ref{alg:discover}--\ref{alg:dispatch} detail these interactions; unattended repair and resubmission are not implemented.

\begin{algorithm}[H]
\caption{DiscoverClusterTopology: archive resource observations}\label{alg:discover}
\begin{algorithmic}[1]
\Require Scheduler read interfaces
\Ensure Timestamped snapshot $G_t$ with explicit unknowns
\State $R\gets\emptyset$; $t\gets\Call{CurrentTime}{}$
\For{$q$ in \{node inventory, partitions, topology\}}
\State $R[q]\gets\Call{ReadWithTimeout}{q}$ \Comment{output, status or unavailable}
\EndFor
\State \Return $(R,t,\text{bandwidth unknown},\text{latency unknown})$
\end{algorithmic}
\end{algorithm}
\begin{algorithm}[H]
\caption{ExecuteSmokeTestLoop: assess a bounded iteration}\label{alg:smoke}
\begin{algorithmic}[1]
\Require Plan $P$, accumulated evidence $E$, snapshot $G_t$
\Ensure READY, REVALIDATE with proposal, or BLOCKED with reasons
\State \Call{CheckRequestFeasibility}{$P,G_t$} \Comment{operator checks site constraints}
\State $r\gets\Call{AssessReadiness}{P,E}$ \Comment{Equation~\ref{eq:ready}}
\If{$r.\mathrm{ready}$} \State \Return READY, $r$ \EndIf
\If{attempt or compute budget exhausted} \State \Return BLOCKED, $r$ \EndIf
\State $P'\gets\Call{ProposeRefinement}{P,E}$
\If{$P'\neq\bot$}
\State \Return REVALIDATE, $P'$ \Comment{collect new matching evidence}
\EndIf
\State \Return BLOCKED, $r$ \Comment{diagnose fault or obtain missing probes}
\end{algorithmic}
\end{algorithm}
\begin{algorithm}[H]
\caption{OrchestrateFullExperiment: operator-facing dispatch}\label{alg:dispatch}
\begin{algorithmic}[1]
\Require Plan $P$, evidence $E$, dispatch mode $d$ (dry-run by default)
\State $r\gets\Call{AssessReadiness}{P,E}$
\State $r.\mathrm{reasons}\gets r.\mathrm{reasons}\cup\Call{VerifyDeclaredInputs}{P}$
\State $r.\mathrm{reasons}\gets r.\mathrm{reasons}\cup\Call{ValidateExecutionGraph}{P}$
\State \Call{PrintDecision}{r}
\If{$r.\mathrm{reasons}\neq\emptyset$} \State \Return BLOCKED \EndIf
\State \Return \Call{UpstreamSubmit}{d} \Comment{explicit operator submission}
\end{algorithmic}
\end{algorithm}

\subsection{Resource discovery and request selection}
Algorithm~\ref{alg:discover} performs light metadata reads. Each read is bounded in time and preserves failures, preventing unavailable topology from silently becoming an apparently healthy empty cluster. Discovery archives what the operator can inspect; it does not place jobs.

Request selection respects authorized accounts, partitions, accelerator type, memory and walltime. Let $\mathcal{P}(j,t)$ denote the scheduler's priority for job $j$ at time $t$. It is observed or queried where available; it is not a value that \tool{} can independently increase. Selection resides in reviewed configuration. Short requests may create backfill opportunities, but their effect on waiting time must be measured. Discovery and monitoring must respect site limits on scheduler queries.

\subsection{Coverage-directed validation}
An obligation list explains why a probe is needed before execution. In the running example, imports and initialization test deployment, a boundary input tests a declared envelope, an output check tests schema/cardinality, a checkpoint probe tests recovery, and a collective probe tests communication. A test cannot satisfy an unexercised obligation by completing quickly. A single-GPU observation must not be relabeled as evidence for a distributed workload.

Algorithm~\ref{alg:smoke} uses the two most recent observations for the exact plan. Older or different-plan records do not establish readiness but retain their compute charges. A recent failed observation blocks admission even if an earlier attempt passed. Two observations reduce dependence on one accidental success; they do not characterize rare faults. Evidence expiry is a policy choice, not an empirically discovered boundary of resource stability.

The lifetime, memory envelope and validation budget are predeclared policy parameters rather than universal constants. Their experimental values appear in Section~\ref{sec:experimental-setup}. A production study must justify these choices for its workload and distinguish framework-allocated memory from all memory used on a device.

\subsection{Bounded refinement and scientific preservation}
A feedback mechanism needs permitted actions and stopping rules. The approach permits one action: halve a batch size greater than one after an out-of-memory classification, only if the scientific contract declares batch invariance and the failing evidence matches the plan. The action modifies the execution subdocument and never the scientific one, and it is returned as a proposal for operator review rather than executed.

The revised hash prevents old records from validating the new configuration. Two new matching observations are needed, with previous costs retained. Unknown faults produce no proposal. Attempt and compute caps bound assessment eligibility, and scheduler limits bound the actual allocation. An autonomous wall-clock budget controller is outside the scope of the approach.

Preserving JSON fields is necessary but insufficient for preserving the scientific estimand. Training batch changes can alter optimization; inference scheduling and numerical nondeterminism can alter outputs. The prospective campaign therefore tests permitted repairs against predeclared tolerances. A model substitution, precision change or revised dataset creates a new scientific specification rather than an invisible recovery action.

\subsection{Guarded execution and agent access}
The execution backend expands the experiment graph, submits dependencies, runs probes and retains provenance. These mechanisms are inherited from \upstream{}; the additional contribution is the admission contract connecting a particular plan to evidence and a bounded refinement proposal.

Algorithm~\ref{alg:dispatch} reassesses readiness immediately before dispatch and checks whether declared inputs and the resolved execution graph still match the assessed plan. An unsupported mapping or changed input blocks the request. A successful check produces a dry run unless the operator explicitly requests submission. Scientific metadata declares intended behavior; it does not prove the scientific semantics of arbitrary commands. The gate is a workflow convention, not a scheduler security boundary.

Agents access discovery, assessment and refinement proposals through a restricted MCP interface, while an operator controls evidence collection and dispatch. This separation makes the decision and its reasons inspectable without granting an agent unrestricted execution capability. Section~\ref{sec:implementation} describes how each component is realized, including modules, protocol version, supported mappings and implementation limits.

\section{Related Work}\label{sec:related}
\subsection{Workflow execution and provenance}
Pegasus, Snakemake, Nextflow and Parsl address dependencies, resource execution and portability \cite{deelman2015pegasus,molder2021snakemake,ditommaso2017nextflow,babuji2019parsl}. CWL standardizes reuse; Sumatra and Sacred capture experimental context \cite{crusoe2022cwl,davison2012sumatra,greff2017sacred}. An auditable workflow and environment capture are therefore foundations, not novel \tool{} features. A container improves mobility but does not itself establish probe coverage or freshness after a resource change.

Our contribution is a particular admission contract above execution machinery: mandatory evidence must still apply to the exact specification, and its acquisition cost must remain bounded. Such a policy can complement established workflow systems; this artifact integrates only the supplied Slurm template. We do not claim other systems cannot implement equivalent policies.

\subsection{Agents that execute and repair workflows}
SWE-agent and OpenHands use agent-computer interfaces for substantial software tasks \cite{yang2024sweagent,wang2025openhands}. ExecutionAgent executes tests of arbitrary projects, making iterative environment construction a relevant comparator \cite{bouzenia2025executionagent}. MLAgentBench evaluates agents on machine-learning experimentation \cite{huang2024mlagentbench}. These studies make feedback-driven execution central and show why static-only comparison cannot isolate the value of \tool{}.

Setup recovery differs from readiness for an expensive production allocation. Our protocol retains native recovery and checkpoints and includes smoke-then-production as a strong baseline. MCP standardizes tools but does not establish their observations' validity. The contribution is exposure of a bounded assessment contract, not a claim that a standardized interface makes experiments reliable.

\subsection{Empirical rigor and artifact evaluation}
Reproducibility work emphasizes methods and execution context; empirical standards make study design and evidence explicit \cite{pineau2021improving,ralph2021standards}. EMSE case-study guidance emphasizes defined cases, collection procedures and threats \cite{runeson2009guidelines}. Scientific benchmarking requires measurement variability and execution conditions to be considered \cite{hoefler2015scientific}. Together, these motivate separating controller conformance, integration feasibility and field effectiveness.
\begin{table}[t]
\caption{Relationship to established mechanisms. Entries describe relevant roles, not exhaustive feature comparisons.}\label{tab:position}
\small\begin{tabularx}{\textwidth}{lXX}
\toprule Mechanism & Established role & Additional readiness question \\
\midrule
Workflow engine & Dependencies and execution & Does evidence authorize this plan? \\
Context/container & Identity and portability & Were mandatory probes passed? \\
Agent setup repair & Iterative environment construction & Are cost and scientific changes bounded? \\
Slurm/backfill & Policy-controlled allocation & Is the request feasible and justified? \\
\upstream{} templates & Smoke runs and audit & Is evidence current, complete and stable? \\
\bottomrule\end{tabularx}
\end{table}
